This has been one of the most crucial points of the whole work.
First of all I have tried to instrument the forkserver function directly inline.
This approach present many problems. The biggest difficulty is the fact that a forkserver needs some library functions 
to work properly :
\begin{itemize}
    \item \texttt{read/write} to implement the communication protocol with the fuzzer.
    \item \texttt{fork/waitpid} to implement the forkserver main logic.
\end{itemize}
These functions could be part of the target binary (if the binary already uses these functions) or not.
In the best case the binary already uses these functions and they are available. 
In this case a call to these functions could be instrumented at any time taking into consideration the type of 
linking with which the binary has been comppiled.
\begin{itemize}
    \item In case of \textit{Static Linking} the function instructions are already inside the target application, 
        because during the compilation the linker inserted it. So a call to these functions could be inserted finding 
        the function address and then instrument a call instruction to that address.
    \item In case of \textit{Dynamic Linking} the function instructions are not inside the target application when 
        the application is launched, but the function addresses are dynamically resolved by the dynamic linker at need. 
        In this case two new section must be used (\texttt{.plt} and \texttt{.got} sections as seen in 
        \cref{sec:loading_relocations}).
\end{itemize}

If the functions required by the forkserver are part of the target binary, their calls can be easily instrumented. 
Depending on whether the binary is statically or dynamically linked, this can be done either by finding the address 
of the function and calling it directly, or by finding the address of its .plt entry and calling it, 
leaving the dynamic linker to resolve the function. 
The problem arises when the original binary does not include these functions. 
In this case, it is almost impossible to modify the .plt and .got sections to include the new functions 
while preserving the binary layout and correctness in the general case. In fact, in some exceptional cases, 
padding may provide enough space in the .got section for a new entry, but this cannot be assumed in the general case. 
Therefore, this first approach has been completely discarded.
Since the forkserver is more complicated than the simple shared map update due to the presence of many controls on 
environment variables and communication with the fuzzer, 
the best approach is to inject the compiled forkserver code inside the target through a \textit{Shared Onbect Injection}.
In order to achieve a Shared Object Injection the \textit{.dynamic} section of the target binary must be modified.
In particular :
\begin{itemize}
    \item A \textit{DT\_NEEDED} entry is added with the name of the shared object to be injected. \footnote{ In reality 
        the shared object name is a string that is written into the \textit{.dynstr} section and into the 
        \textit{.dynamic} section is specified the offset of the string written into \textit{.dynstr}}
    \item A \textit{RUNPATH} entry is added containing the path in which is saved the actual shared object.
\end{itemize}
\footnote{
    The injection could be also made with LD\_PRELOAD environment variable. 
    This variable tells the dynamic linker to add the specified libraries as dependencie of the executable.
    The problem is that at each execution the variable must be resetted with the name of the shared object.
    Adding the DT\_NEEDED dependency instead the dependency is persistent
}\\
The shared object injected into the target binary, contains the code of the forkserver that must be executed 
prior the main of the binary.
In order to achieve the fact that a function of a shared object injected into a binary is executed before the main, 
that function must be declared as a ``constructor''.
As stated in \cref{sec:loading_relocations}, by declaring a function as constructor, the address of that function is 
placed inside the \textit{.init\_array} section and executed before the actual start of the main of the executable.
When an executable has some dependencies, the dynamic linker execute the \textit{.init} section and the functions 
contained into \textit{.init\_array} section before actually start the main of the executable.\footnote{
    The .init section is the legacy version of .init\_array section. 
}
In \cref{fig:dl_sub2} can be seen the execution flow of the dynamic linker when it initialize the dynamic dependencies of a target binary that 
has the dynamic dependencies shown in \cref{fig:dl_sub1}.
\begin{enumerate}
    \item The dynamic linker search for \textit{DT\_NEEDED} entries of the target binary. When a dependency is found, it map the library 
        into a loaded \gls{VMA} of the process associated with the target binary that is being executed.
    \item For each library founded previously, the dynamic linker execute each function found in the \textit{.init\_array} 
        section of that shared object.
    \item When the dynamic linker ends, the \texttt{\_start} function is executed and it execute \texttt{\_\_libc\_start\_main} 
        function that execute all the functions in the \textit{.init\_array} of the target binary.
\end{enumerate}
\begin{figure}[ht]
\centering
        \begin{subfigure}[h]{0.7\textwidth}
        \centering
        \includegraphics[width=\linewidth]{init_array_loading_example_a.png}
        \caption{Layout of a target binary with its dynamic dependencies}
        \label{fig:dl_sub1}
    \end{subfigure}
    \vspace{0.5cm}
    \begin{subfigure}[h]{0.7\textwidth}
        \centering
        \includegraphics[width=\linewidth]{init_array_loading_example_b.png}
        \caption{Dynamic linker initialization of dynamic dependencies}
        \label{fig:dl_sub2}
    \end{subfigure}
    \caption{An example of ELF execution with dynamic dependencies and constructor function}
    \label{fig:constructor_function_execution}
\end{figure}
Therefore, injecting a \gls{SO} that has a function declared as constructor into the target binary, could resolve the 
problem of instrument the forkserver in such a way that it is executeb prior the main.
